\documentclass[11pt,a4paper]{article}

\usepackage[a4paper,margin=1.5cm]{geometry}
\usepackage[T1]{fontenc}
\usepackage[utf8]{inputenc}
\usepackage{lmodern}
\usepackage{amsmath,amssymb}
\usepackage{array}
\usepackage{enumitem}
\usepackage{fancyhdr}
\usepackage{hyperref}
\usepackage{longtable}
\usepackage{tabularx}
\usepackage{xcolor}

\hypersetup{
  colorlinks=true,
  linkcolor=black,
  urlcolor=blue
}

\pagestyle{fancy}
\fancyhf{}
\lhead{\textsc{BCS2410 Mock Examination Pack}}
\rhead{\textsc{Page \thepage}}
\cfoot{}

\setlength{\parindent}{0pt}
\setlength{\parskip}{0.45em}
\setlength{\headheight}{14pt}
\setlist[enumerate]{leftmargin=*,itemsep=0.7em}

\newcommand{\paperheader}[3]{
  \begin{center}
  {\Large\textbf{BCS2410 Embedded Programming}}\\[0.25em]
  {\large\textsc{Mock Examination Paper}}\\[0.35em]
  {\Large\textbf{#1}}
  \end{center}
  \vspace{0.25em}
  \begin{tabularx}{\linewidth}{|>{\raggedright\arraybackslash}p{0.22\linewidth}|X|}
  \hline
  Total marks & #2 \\
  \hline
  Recommended duration & #3 \\
  \hline
  Paper structure & Section A = 40 marks of short theory / closed-ended questions. Section B = 60 marks of practical exercises. \\
  \hline
  Candidate instructions & Answer all questions. Show working where appropriate. Write assembly clearly with labels, branch conditions, and register use visible. Where a question requests explanation, concise technical prose is sufficient. \\
  \hline
  \end{tabularx}
  \vspace{0.6em}
  \begin{tabularx}{\linewidth}{|p{0.23\linewidth}|X|p{0.22\linewidth}|X|}
  \hline
  Candidate name & & Student number & \\
  \hline
  Date & & Start time & \\
  \hline
  \end{tabularx}
  \vspace{0.75em}
}

\newcommand{\subpaperheader}[1]{
  \begin{center}
  \fbox{
  \begin{minipage}{0.97\linewidth}
  \textbf{#1}
  \end{minipage}}
  \end{center}
}

\newcommand{\qmarks}[1]{\hfill{\textbf{[#1 marks]}}}
\newcount\answerlinecount

\newcommand{\answerlines}[1]{
  \par\smallskip\noindent\textit{Answer:}\par\nobreak
  \answerlinecount=0
  \loop\ifnum\answerlinecount<#1
    \par\noindent\rule{\linewidth}{0.4pt}\vspace{0.95\baselineskip}
    \advance\answerlinecount by1
  \repeat
  \par\smallskip
}

\newcommand{\tinyanswer}{\answerlines{2}}
\newcommand{\smallanswer}{\answerlines{4}}
\newcommand{\medanswer}{\answerlines{6}}
\newcommand{\largeanswer}{\answerlines{9}}
\newcommand{\xlargeanswer}{\answerlines{13}}

\newcommand{\checkbox}{$\square$}

\begin{document}
\pagenumbering{roman}

\begin{titlepage}
\centering
{\Large \textsc{University Examination Style Revision Pack}}\\[1.0em]
{\Huge \textbf{BCS2410 Embedded Programming}}\\[0.4em]
{\LARGE \textbf{Mock Examination Pack}}\\[1.5em]

{\large Four simulated exam papers with answer space}\\[0.5em]
{\large Built from the lecture notes, cram sheets, self-tests, worked drills, and syllabus-aligned exam notes in this vault}\\[2em]

\begin{center}
\fbox{
\begin{minipage}{0.92\textwidth}
\textbf{How this pack is designed}
\begin{itemize}[leftmargin=1.2em]
\item Each paper is worth \textbf{100 marks}.
\item Each paper follows the exam structure inferred from your notes: about \textbf{40\% theory / short closed questions} and \textbf{60\% practical work in 2 or 3 larger exercises}.
\item Across the four papers, all major course areas are covered: foundations, architecture, data representation, C and pointers, ARM Thumb assembly, peripherals and workflow, interrupts and clocks, and FPGA / Edge AI.
\item The difficulty is intentionally mixed so you are comfortable whether the real exam is straightforward or demanding.
\item A \textbf{condensed model solution and marking guide} is included at the end of this PDF.
\end{itemize}
\end{minipage}}
\end{center}

\vfill

\begin{center}
\fbox{
\begin{minipage}{0.92\textwidth}
\textbf{Exam-use suggestion}
\begin{itemize}[leftmargin=1.2em]
\item First attempt a paper under timed conditions.
\item Mark it only after you have completed the full paper.
\item Use the answer boxes as if this were the actual answer script.
\end{itemize}
\end{minipage}}
\end{center}

\vfill

Generated in the notes vault: \\
\url{/Users/davidwickerhf/Projects/personal/notes/University/BCS2410 Embedded Programming Knowledge Base/12 Mock Exams/}
\end{titlepage}

\clearpage
\tableofcontents
\clearpage
\pagenumbering{arabic}

\section*{Paper 1}
\addcontentsline{toc}{section}{Paper 1}
\paperheader{Paper 1: Core Embedded Systems, Data Representation, and ARM Basics}{100}{2 hours 10 minutes}

\subpaperheader{Section A: Theory and Short Closed Questions (40 marks)}
\begin{enumerate}
\item State one key difference between an embedded system and a general-purpose computer. \qmarks{4}
\tinyanswer

\item Program code on the Discovery board is typically stored in:
\begin{itemize}[leftmargin=1.5em]
\item \checkbox\ Random Access Memory (RAM)
\item \checkbox\ Flash / ROM-like non-volatile memory
\item \checkbox\ Cache only
\item \checkbox\ The Arithmetic Logic Unit (ALU)
\end{itemize}
\tinyanswer

\item What does the Program Counter (PC) hold? \qmarks{4}
\tinyanswer

\item Name the two main parts of the processor in the datapath-plus-control view. \qmarks{4}
\tinyanswer

\item In a little-endian system, how is \texttt{0x12345678} stored in memory from lowest to highest address? \qmarks{4}
\tinyanswer

\item What is the main job of the Arithmetic Logic Unit (ALU)? \qmarks{4}
\tinyanswer

\item Match each flag to its meaning: \texttt{N}, \texttt{Z}, \texttt{C}, \texttt{V}. \qmarks{4}
\smallanswer

\item What is the difference between \texttt{MOV}, \texttt{LDR}, and \texttt{STR}? \qmarks{4}
\smallanswer

\item Which branch is appropriate for a signed greater-than comparison after \texttt{CMP}: \texttt{BGT} or \texttt{BHI}? Briefly justify. \qmarks{4}
\smallanswer

\item What does \texttt{*(ptr + 2)} mean when \texttt{ptr} is an \texttt{int *} pointing at the first element of an array? \qmarks{4}
\smallanswer
\end{enumerate}

\subpaperheader{Section B: Practical Exercises (60 marks)}

\textbf{Exercise 1: ARM Thumb loop and branch construction} \qmarks{20}

Write an ARM Thumb function \texttt{sum\_to\_n} with the following behaviour:
\begin{itemize}[leftmargin=1.5em]
\item input: \texttt{r0 = n}
\item output: return \texttt{1 + 2 + \dots + n} in \texttt{r0}
\item if \texttt{n <= 0}, return \texttt{0}
\end{itemize}

Your answer should:
\begin{enumerate}[label=\alph*)]
\item use a loop with labels \qmarks{10}
\item use the correct signed branch conditions \qmarks{5}
\item briefly explain why \texttt{BGT} and not \texttt{BHI} is the right exit-test style here \qmarks{5}
\end{enumerate}
\xlargeanswer

\textbf{Exercise 2: Pointers, arrays, and memory layout} \qmarks{20}

Consider the following C code:
\begin{verbatim}
int arr[] = {10, 20, 30, 40};
int *ptr = arr;
int a = *(ptr + 2);
*(ptr + 1) = 99;
\end{verbatim}

Answer all parts:
\begin{enumerate}[label=\alph*)]
\item What does \texttt{ptr} store? \qmarks{4}
\item What is the value of \texttt{a}? \qmarks{4}
\item What are the final contents of \texttt{arr}? \qmarks{4}
\item Explain why \texttt{ptr + 2} means ``two elements later'' and not ``two bytes later.'' \qmarks{4}
\item Show the byte layout of \texttt{0x12345678} in little-endian memory. \qmarks{4}
\end{enumerate}
\xlargeanswer

\textbf{Exercise 3: Memory-mapped I/O and load-modify-store} \qmarks{20}

A General-Purpose Input/Output (GPIO) output register address is already stored in \texttt{r0}. Write ARM Thumb assembly that:
\begin{enumerate}[label=\alph*)]
\item reads the register \qmarks{4}
\item sets bit 7 to 1 \qmarks{4}
\item clears bit 3 to 0 \qmarks{4}
\item writes the result back \qmarks{4}
\item briefly explains why ordinary \texttt{LDR} and \texttt{STR} can be used for peripheral registers \qmarks{4}
\end{enumerate}
\largeanswer
\largeanswer

\newpage
\section*{Paper 2}
\addcontentsline{toc}{section}{Paper 2}
\paperheader{Paper 2: Architecture, Flags, Functions, and Assembly Tracing}{100}{2 hours 10 minutes}

\subpaperheader{Section A: Theory and Short Closed Questions (40 marks)}
\begin{enumerate}
\item What is the difference between datapath and control? \qmarks{4}
\tinyanswer

\item What is the role of the decode stage in the fetch-decode-execute cycle? \qmarks{4}
\tinyanswer

\item Why are registers faster than ordinary memory access? \qmarks{4}
\tinyanswer

\item State one difference between Von Neumann and Harvard architecture. \qmarks{4}
\tinyanswer

\item What is a half adder, and what extra input makes a full adder different? \qmarks{4}
\smallanswer

\item What does \texttt{BL} do that \texttt{B} does not? \qmarks{4}
\tinyanswer

\item Where are the first four 32-bit function arguments usually passed in ARM Cortex-M style calling convention? \qmarks{4}
\tinyanswer

\item What is the purpose of the Link Register (\texttt{LR})? \qmarks{4}
\tinyanswer

\item Why can \texttt{BGT} and \texttt{BHI} give different answers on the same raw bit patterns? \qmarks{4}
\smallanswer

\item What is cache for? \qmarks{4}
\tinyanswer
\end{enumerate}

\subpaperheader{Section B: Practical Exercises (60 marks)}

\textbf{Exercise 1: Function calls, \texttt{BL}, \texttt{LR}, and tracing} \qmarks{30}

Consider the following ARM Thumb code:
\begin{verbatim}
ASM_Function:
    MOVS r0, #3
    BL   add5
    BL   twice
    BX   LR

add5:
    ADDS r0, r0, #5
    BX   LR

twice:
    LSL  r0, r0, #1
    BX   LR
\end{verbatim}

Answer all parts:
\begin{enumerate}[label=\alph*)]
\item Trace the value in \texttt{r0} after each instruction in \texttt{ASM\_Function}. \qmarks{8}
\item What final value is returned in \texttt{r0}? \qmarks{4}
\item What does each \texttt{BL} instruction do to control flow and to \texttt{LR}? \qmarks{8}
\item Explain why \texttt{BX LR} returns correctly in this example. \qmarks{4}
\item If \texttt{add5} itself needed to call another function, why might it need to preserve \texttt{LR} first? \qmarks{6}
\end{enumerate}
\xlargeanswer
\largeanswer

\textbf{Exercise 2: Flags, signedness, overflow, and bitwise reasoning} \qmarks{30}

Answer all parts.

\begin{enumerate}[label=\alph*)]
\item Given:
\begin{verbatim}
MOV  r0, #-1
MOV  r1, #2
CMP  r0, r1
\end{verbatim}
State whether each branch would be taken, and explain why:
\texttt{BGT}, \texttt{BLT}, \texttt{BHI}, \texttt{BLO}. \qmarks{12}

\item Compute the result of:
\begin{verbatim}
10110100 AND 00001111
10110100 OR  00001111
10110100 XOR 00001111
\end{verbatim}
\qmarks{6}

\item A 32-bit signed integer currently holds \texttt{0x7FFFFFFF}. What happens when 1 is added? State the new bit pattern and explain the signed interpretation. \qmarks{6}

\item Write a short ARM Thumb sequence that sets \texttt{r1 = 1} if \texttt{r0} is odd, else \texttt{r1 = 0}. You may use \texttt{TST} or a branch-based solution. \qmarks{6}
\end{enumerate}
\xlargeanswer
\xlargeanswer

\newpage
\section*{Paper 3}
\addcontentsline{toc}{section}{Paper 3}
\paperheader{Paper 3: Workflow, Debugging, Peripherals, Clocks, and Integration}{100}{2 hours 10 minutes}

\subpaperheader{Section A: Theory and Short Closed Questions (40 marks)}
\begin{enumerate}
\item What is the difference between build, flash, and debug? \qmarks{4}
\smallanswer

\item What is ST-LINK used for? \qmarks{4}
\tinyanswer

\item Why is the disassembly view useful during debugging? \qmarks{4}
\smallanswer

\item What is a breakpoint? \qmarks{4}
\tinyanswer

\item What is memory-mapped I/O? \qmarks{4}
\smallanswer

\item What is the difference between a datasheet and a programming manual? \qmarks{4}
\smallanswer

\item What is an interrupt? \qmarks{4}
\tinyanswer

\item What is an Interrupt Service Routine (ISR)? \qmarks{4}
\tinyanswer

\item Why do peripherals need clocks? \qmarks{4}
\tinyanswer

\item State one advantage and one disadvantage of polling. \qmarks{4}
\smallanswer
\end{enumerate}

\subpaperheader{Section B: Practical Exercises (60 marks)}

\textbf{Exercise 1: Debugging workflow on the Discovery board} \qmarks{20}

You build a project successfully, but the board does not behave as expected and an LED never changes state.

Describe a structured debug workflow using course tools and concepts. Your answer should include:
\begin{enumerate}[label=\alph*)]
\item what to check after build \qmarks{4}
\item what flashing does \qmarks{4}
\item how ST-LINK is involved \qmarks{4}
\item how breakpoints, register views, memory views, and disassembly can help isolate the bug \qmarks{8}
\end{enumerate}
\xlargeanswer

\textbf{Exercise 2: Integrating ARM assembly with C} \qmarks{20}

You want to call an assembly function from C with this prototype:
\begin{verbatim}
extern unsigned max3u(unsigned a, unsigned b, unsigned c);
\end{verbatim}

Write the assembly implementation for \texttt{max3u} that returns the maximum of the three \textbf{unsigned} inputs.

Requirements:
\begin{enumerate}[label=\alph*)]
\item use argument registers correctly \qmarks{4}
\item use unsigned comparison branches correctly \qmarks{8}
\item return the result in the correct register \qmarks{4}
\item briefly explain why signed branches would be wrong here \qmarks{4}
\end{enumerate}
\xlargeanswer

\textbf{Exercise 3: Polling, interrupts, timers, and control-flow design} \qmarks{20}

You are asked to design a simple embedded application with this behaviour:
\begin{itemize}[leftmargin=1.5em]
\item an LED toggles every 500 ms
\item a button press must be detected promptly
\item the CPU should not waste time unnecessarily
\end{itemize}

Answer all parts:
\begin{enumerate}[label=\alph*)]
\item Would you prefer polling, interrupts, or a mixture of both? Justify your choice. \qmarks{8}
\item Explain the role of a timer in this design. \qmarks{4}
\item Sketch a high-level control flow for the solution. Pseudocode is acceptable. \qmarks{8}
\end{enumerate}
\xlargeanswer
\largeanswer

\newpage
\section*{Paper 4}
\addcontentsline{toc}{section}{Paper 4}
\paperheader{Paper 4: FPGA, Edge AI, and Integrated Embedded System Reasoning}{100}{2 hours 10 minutes}

\subpaperheader{Section A: Theory and Short Closed Questions (40 marks)}
\begin{enumerate}
\item What is a Field-Programmable Gate Array (FPGA)? \qmarks{4}
\tinyanswer

\item State one key difference between an FPGA and an Application-Specific Integrated Circuit (ASIC). \qmarks{4}
\tinyanswer

\item What is a Lookup Table (LUT) used for in FPGA fabric? \qmarks{4}
\tinyanswer

\item What is the role of Block RAM (BRAM)? \qmarks{4}
\tinyanswer

\item Why do modern FPGA platforms include Digital Signal Processing (DSP) blocks? \qmarks{4}
\tinyanswer

\item What is a Deep Processing Unit (DPU) in the course context? \qmarks{4}
\tinyanswer

\item What is quantization, and why is it useful on edge devices? \qmarks{4}
\smallanswer

\item What is a bitstream? \qmarks{4}
\tinyanswer

\item In a Zynq-style platform, what is the difference between the processing system and programmable logic? \qmarks{4}
\smallanswer

\item What does Vitis AI broadly do in the deployment flow? \qmarks{4}
\smallanswer
\end{enumerate}

\subpaperheader{Section B: Practical Exercises (60 marks)}

\textbf{Exercise 1: Edge AI deployment scenario} \qmarks{30}

You must deploy a small image-classification model on an embedded platform with limited power and memory.

Answer all parts:
\begin{enumerate}[label=\alph*)]
\item Compare three implementation options: CPU only, FPGA with DPU, and ASIC. Focus on flexibility, development cost, and runtime efficiency. \qmarks{12}
\item Explain a sensible Vitis AI style workflow from model preparation to execution on hardware. \qmarks{8}
\item Explain where quantization fits into this workflow and why it is often helpful. \qmarks{5}
\item State the likely roles of LUTs, BRAM, and DSP blocks in this kind of design. \qmarks{5}
\end{enumerate}
\xlargeanswer
\largeanswer

\textbf{Exercise 2: Mixed embedded-system practical design} \qmarks{30}

Consider a smart sensor node with:
\begin{itemize}[leftmargin=1.5em]
\item a microcontroller
\item Flash and Random Access Memory (RAM)
\item a timer
\item General-Purpose Input/Output (GPIO)
\item a status register at address \texttt{0x40001000}
\item a data register at address \texttt{0x40001004}
\end{itemize}

Bit 0 of the status register becomes 1 when new sensor data is ready.

Answer all parts:
\begin{enumerate}[label=\alph*)]
\item State where program code should live and where live runtime variables should live. \qmarks{6}
\item Explain why GPIO, timers, and sensor registers count as peripherals rather than the CPU itself. \qmarks{6}
\item Write ARM Thumb assembly that repeatedly polls the status register until bit 0 is set, then loads the data register value into \texttt{r2}. \qmarks{12}
\item Briefly explain why this is an example of memory-mapped I/O. \qmarks{6}
\end{enumerate}
\xlargeanswer
\largeanswer

\newpage
\section*{Condensed Model Solutions and Marking Guide}
\addcontentsline{toc}{section}{Condensed Model Solutions and Marking Guide}

\begin{center}
\fbox{
\begin{minipage}{0.97\linewidth}
Use this section only after attempting a paper. Solutions are intentionally concise and aligned to the vault's course notes rather than to every possible alternative answer.
\end{minipage}}
\end{center}

\subsection*{Paper 1 Solutions}
\textbf{Section A}
\begin{enumerate}
\item Embedded systems are specialized, hardware-coupled, and resource-constrained; general-purpose systems are meant for broad use.
\item Flash / ROM-like non-volatile memory.
\item The address of the next instruction to execute.
\item Datapath and control.
\item \texttt{78 56 34 12}.
\item Arithmetic, logic, comparisons, shifts, and related value transformations.
\item \texttt{N} negative, \texttt{Z} zero, \texttt{C} carry / no-borrow style unsigned condition, \texttt{V} signed overflow.
\item \texttt{MOV} moves register / immediate values; \texttt{LDR} loads from memory to a register; \texttt{STR} stores a register to memory.
\item \texttt{BGT}. It is a signed comparison; \texttt{BHI} is for unsigned higher.
\item The value two integer elements after \texttt{ptr}; if \texttt{ptr} starts at the first element, this is the third element.
\end{enumerate}

\textbf{Section B}

\textit{Exercise 1 reference solution}
\begin{verbatim}
sum_to_n:
    CMP  r0, #0
    BLE  done_zero
    MOVS r1, #1
    MOVS r2, #0
loop:
    CMP  r1, r0
    BGT  done_sum
    ADD  r2, r2, r1
    ADDS r1, r1, #1
    B    loop
done_sum:
    MOV  r0, r2
    BX   LR
done_zero:
    MOVS r0, #0
    BX   LR
\end{verbatim}
Marks: correct structure 10, correct signed branches 5, signed-vs-unsigned explanation 5.

\textit{Exercise 2 key points}
\begin{itemize}[leftmargin=1.5em]
\item \texttt{ptr} stores the address of \texttt{arr[0]}.
\item \texttt{a = 30}.
\item Final array: \texttt{\{10, 99, 30, 40\}}.
\item Pointer arithmetic scales by element size.
\item Little-endian layout of \texttt{0x12345678}: \texttt{78 56 34 12}.
\end{itemize}

\textit{Exercise 3 reference solution}
\begin{verbatim}
    LDR  r1, [r0]
    ORR  r1, r1, #0x80
    BIC  r1, r1, #0x08
    STR  r1, [r0]
\end{verbatim}
\texttt{BIC} clears the selected bit. Full credit also for a correct clear-by-mask alternative. Peripheral registers occupy normal addresses, so \texttt{LDR} and \texttt{STR} work.

\subsection*{Paper 2 Solutions}
\textbf{Section A}
\begin{enumerate}
\item Datapath moves / transforms values; control issues signals that steer execution.
\item Decode interprets the fetched instruction and generates the internal control actions.
\item Registers are inside the CPU and much closer to the Arithmetic Logic Unit.
\item Von Neumann shares instruction and data memory; Harvard separates them.
\item A half adder adds two bits; a full adder also accepts carry-in.
\item \texttt{BL} branches and stores a return address in \texttt{LR}.
\item In \texttt{r0-r3}.
\item It holds the return address after a call-style branch.
\item Signed and unsigned branches interpret the same flag state differently.
\item It reduces average access time by holding frequently used data / instructions close to the CPU.
\end{enumerate}

\textbf{Section B}

\textit{Exercise 1}
\begin{itemize}[leftmargin=1.5em]
\item \texttt{MOVS r0, \#3} gives \texttt{r0 = 3}.
\item After \texttt{BL add5}, \texttt{r0 = 8}.
\item After \texttt{BL twice}, \texttt{r0 = 16}.
\item Final returned value is \texttt{16}.
\item Each \texttt{BL} changes control flow and writes the return address into \texttt{LR}.
\item \texttt{BX LR} works because the callee returns to the address saved by \texttt{BL}.
\item If \texttt{add5} made another call, that new \texttt{BL} would overwrite \texttt{LR}; it may need \texttt{PUSH \{lr\}} / \texttt{POP \{lr\}} or equivalent preservation.
\end{itemize}

\textit{Exercise 2}
\begin{itemize}[leftmargin=1.5em]
\item With \texttt{r0 = -1} and \texttt{r1 = 2}: \texttt{BGT} no, \texttt{BLT} yes, \texttt{BHI} yes, \texttt{BLO} no.
\item \texttt{10110100 AND 00001111 = 00000100}
\item \texttt{10110100 OR 00001111 = 10111111}
\item \texttt{10110100 XOR 00001111 = 10111011}
\item \texttt{0x7FFFFFFF + 1 = 0x80000000}; under signed 32-bit interpretation this is negative and represents signed overflow.
\item One correct odd-test sequence:
\begin{verbatim}
    TST  r0, #1
    BEQ  even_case
    MOVS r1, #1
    B    done
even_case:
    MOVS r1, #0
done:
\end{verbatim}
\end{itemize}

\subsection*{Paper 3 Solutions}
\textbf{Section A}
\begin{enumerate}
\item Build creates the binary, flash programs it into the microcontroller, debug runs with inspection and breakpoints.
\item ST-LINK is the host-to-board programming and debugging interface.
\item Disassembly shows generated instructions and addresses, linking source behaviour to machine execution.
\item A breakpoint stops execution at a chosen location.
\item Peripherals occupy addresses in the normal processor address space.
\item Datasheet: device capabilities / characteristics; programming manual: architecture, programming model, register semantics.
\item An interrupt redirects normal execution so an event can be handled promptly.
\item The function / routine that runs in response to an interrupt.
\item Peripherals depend on clock sources to operate correctly.
\item Advantage: simple to understand. Disadvantage: wastes CPU time and can increase latency.
\end{enumerate}

\textbf{Section B}

\textit{Exercise 1 expected workflow}
\begin{itemize}[leftmargin=1.5em]
\item Confirm successful build and the correct project image.
\item Flash the image into device non-volatile memory.
\item Use ST-LINK for programming / debug connection.
\item Start a debug session, place breakpoints, single-step, inspect variables, inspect registers, inspect memory, compare with disassembly, and identify where behaviour diverges from expectation.
\end{itemize}

\textit{Exercise 2 reference solution}
\begin{verbatim}
    .global max3u
max3u:
    CMP  r0, r1
    BHS  keep_r0
    MOV  r0, r1
keep_r0:
    CMP  r0, r2
    BHS  done
    MOV  r0, r2
done:
    BX   LR
\end{verbatim}
\texttt{BHS} is unsigned higher-or-same; signed branches would be incorrect because the inputs are explicitly unsigned.

\textit{Exercise 3}
\begin{itemize}[leftmargin=1.5em]
\item Best answer: usually a mixture. Timer interrupt or timer-driven event for periodic LED toggling; interrupt or at least event-based detection for prompt button response; superloop may still exist for background work.
\item The timer provides the 500 ms time base.
\item High-level sketch: initialize GPIO and timer, enable button interrupt or sampled event, in ISR set a flag or toggle state, main loop handles background tasks.
\end{itemize}

\subsection*{Paper 4 Solutions}
\textbf{Section A}
\begin{enumerate}
\item A reconfigurable integrated circuit whose logic can be programmed after manufacturing.
\item FPGA is reconfigurable after manufacture; ASIC is fixed once fabricated.
\item LUTs implement combinational logic.
\item BRAM provides on-chip storage.
\item DSP blocks accelerate arithmetic-heavy operations.
\item A dedicated accelerator in the FPGA fabric for neural-network inference.
\item Quantization reduces numeric precision, usually reducing model size and compute cost for constrained hardware.
\item The configuration data that programs the FPGA fabric.
\item Processing system is the processor side; programmable logic is the configurable hardware fabric.
\item It supports model preparation, quantization, compilation, and runtime execution on supported targets.
\end{enumerate}

\textbf{Section B}

\textit{Exercise 1}
\begin{itemize}[leftmargin=1.5em]
\item CPU only: flexible, easiest to program, but less efficient for inference-heavy workloads.
\item FPGA + DPU: middle ground; configurable and more efficient for targeted acceleration.
\item ASIC: most efficient for one fixed purpose, but inflexible and expensive to design.
\item Workflow: prepare model, quantize, compile for target, generate / use platform bitstream and runtime stack, execute with runtime such as VART.
\item Quantization reduces precision, often to INT8, to lower cost while keeping acceptable accuracy.
\item LUTs implement logic, BRAM stores data and activations, DSP blocks accelerate arithmetic.
\end{itemize}

\textit{Exercise 2}
\begin{itemize}[leftmargin=1.5em]
\item Program code: Flash. Runtime variables / stack / buffers: Random Access Memory.
\item GPIO, timers, and sensor registers are peripherals because they are hardware blocks controlled by the CPU, not the CPU core itself.
\item One correct polling solution:
\begin{verbatim}
    LDR  r0, =0x40001000
poll:
    LDR  r1, [r0]
    TST  r1, #1
    BEQ  poll
    LDR  r1, =0x40001004
    LDR  r2, [r1]
\end{verbatim}
\item This is memory-mapped I/O because the peripheral registers are accessed through normal addresses with normal load instructions.
\end{itemize}

\end{document}
